\section*{第 \ref{ch:g12:past-climates} 章——读取过去的气候}

\begin{solution}{exo:g12:past-climates:1}
(a) \omterm{def:g12:past-climates:proxy}{代用指标}；(b) \omterm{def:g12:past-climates:ice-core}{冰芯}；(c) \omterm{def:g12:past-climates:glacial-cycle}{间冰期}；(d) \omterm{def:g12:past-climates:pollen}{孢粉分析}；
(e) \omterm{def:g12:past-climates:lgm}{末次盛冰期}。
\end{solution}

\begin{solution}{exo:g12:past-climates:2}
(a) 极地冰：当时空气中的二氧化碳和甲烷。
(b) 湖泊沉积物或泥炭：植被，此处为需要温暖
夏季的树木。(c) 海洋沉积物：陆地冰的体积和\omterm{def:g10:layered-ocean:thermocline}{深海}的温度
(d) 树木：逐年变化的生长条件。
\end{solution}

\begin{solution}{exo:g12:past-climates:3}
(a) $-50 - (-120) = \qty{70}{m}$，在 \num{7000} 年内：$\qty{70000}{mm}/\num{7000}
= \qty{10}{mm}$ 每年。(b) $10/3.7 \approx 2.7$：去\omterm{def:g12:past-climates:glacial-cycle}{冰期}的平均上升
速度约为 2006--2018 年的 2.7 倍。
\end{solution}

\begin{solution}{exo:g12:past-climates:4}
一个世纪是十个十年：每世纪 $\qty{2.0}{\degreeCelsius}$。一千
年是一百个十年：每千年 $\qty{20}{\degreeCelsius}$。
\end{solution}

\begin{solution}{exo:g12:past-climates:5}
(a) \num{23000} 年：地轴的岁差；\num{42000} 年：
地轴的倾斜；约 \num{100000} 年：轨道的形状。
(b) $\num{1000000}/\num{41000} \approx 24$ 个循环；$\num{1000000}/\num{100000}
= 10$ 个循环。
\end{solution}

\begin{solution}{exo:g12:past-climates:6}
(a) $266.6 + 7.83 \times (-9.4) \approx 266.6 - 73.6 = \qty{193}{ppm}$：仅
比测得的 \qty{190}{ppm} 高 \qty{3}{ppm}。(b) $266.6 + 7.83 \times 5.5
\approx \qty{310}{ppm}$，比最高测量值高 \qty{11}{ppm}：最暖
点位于直线之下。(c) $\Delta T = (427 - 266.6)/7.83 \approx
\qty{20}{\degreeCelsius}$。该直线仅描述两个量
在自然循环中共同变化的方式，且只限于数据范围之内
（\qty{-11}{\degreeCelsius} 至 \qty{+6}{\degreeCelsius}，174 至
\qty{299}{ppm}）。它不是定律，远超出该范围的数值
不能由此读出。
\end{solution}

\begin{solution}{exo:g12:past-climates:7}
(a) 草本植物，\qty{84}{\percent}：开阔的景观，几乎没有树木，气候寒冷。
(b) 橡树（\qty{42}{\percent}）和榛树（\qty{34}{\percent}）：阔叶
林，夏季温暖。(c) 桦树先扩散，随后是松树；榛树和橡树
最后到达。(d) 树木扩散的速度，仅取决于种子传播和
幼树生长的速度，扩散始于寒冷期中存活下来的地点；
土壤也必须形成。森林可能落后于气候。
\end{solution}

\begin{solution}{exo:g12:past-climates:8}
(a) \qty{300}{cm}，即 \qty{3000}{mm}，历时 \num{8000} 年：$3000/8000 =
\qty{0.375}{mm}$ 每年。(b) $\num{6000} + (470 - 300)/300 \times \num{8000}
\approx \num{6000} + \num{4530} \approx \num{10500}$ 年。(c) $10/0.375
\approx 27$ 年。
\end{solution}

\begin{solution}{exo:g12:past-climates:9}
(a) $7.16 - 4.11 = 3.05$ 百万平方千米，$3.05/4.11 \approx
\qty{74}{\percent}$，即陆地多出的部分。(b) 是的：\omterm{def:g12:past-climates:lgm}{末次盛冰期}时，海平面
低于今天 125 至 \qty{134}{m}，远大于地图上连接婆罗洲与亚洲的
\qty{36}{m}水深。(c) 地图采用今天的海底，此后它并
未完全保持不变（其上有沉积物堆积，地面也
发生了轻微移动）；网格也较粗，每 0.1
度一个点，且低海面位置仅能确定到约十米范围内。
\end{solution}

\begin{solution}{exo:g12:past-climates:10}
(a) \omterm{def:g12:past-climates:lgm}{末次盛冰期}的值 \qty{4.92}{\permil} 较高。
这是可以预料的：当时的大冰盖储存了氧-18 含量低的水，
这些水离开海洋，使其中生长的贝壳富含
氧-18；深层水也更冷。(b) 这样，寒冷的\omterm{def:g12:past-climates:glacial-cycle}{冰期}在图上位于低处，温暖的时期位于高处，
与温度曲线的画法一致。
\end{solution}

\begin{solution}{exo:g12:past-climates:11}
(a) 相关性表明两条记录同升同降，接近一条直线；
它不能说明二者谁导致谁，也不能说明任一方
是原因。(b) 变暖的气候可能使二氧化碳释放进大气；
或者第三种原因，例如地球轨道的变化，可能同时影响两者；
或者二者互为因果。(c) 变化的先后次序，
即哪一个先发生，需对冰及其所含空气进行精确定年；以及
对二氧化碳如何作用于地球辐射的理解
（见\cref{ch:g12:climate-change}）。
\end{solution}

\begin{solution}{exo:g12:past-climates:12}
(a) $\num{1300000}/\num{41000} \approx 32$ 个循环，发生于 260 万至
130 万年前之间；$\num{700000}/\num{98500} \approx 7$ 个循环，覆盖
过去 \num{700000} 年。(b) \num{41000} 年：地轴的倾斜；约
\num{100000} 年：轨道的形状。(c) 一个 \num{41000} 年的循环
被采样约 $\num{41000}/\num{2000} \approx 20$ 次，足以描绘它；
持续一个世纪的变化落在两个样本之间，或落在某一个样本内，
几乎不留痕迹。
\end{solution}

\begin{solution}{exo:g12:past-climates:13}
(a) $192/10 \approx 19$ 与 $198/10 \approx 20$：约快 20 倍。
(b) 与每世纪 \qty{4}{ppm} 相比：$192/4 = 48$ 至 $198/4 \approx 50$
倍；与 \qty{42}{ppm} 相比：$192/42 \approx 4.6$ 至 $198/42 \approx 4.7$
倍。今天的上升速率约为古新世–始新世极热事件的 5 至 50 倍。
(c) 上升量为 $2000 - 900 = \qty{1100}{ppm}$：$1100/42 \approx 26$ 个世纪，
约 \num{2600} 年（以该事件的最快速率计）；$1100/195 \approx 5.6$
个世纪，不足 600 年（以今天的速率计）。
\end{solution}

\begin{solution}{exo:g12:past-climates:14}
(a) 十倍于 1 至 \qty{1.5}{\degreeCelsius}：每千年 10 至 \qty{15}{\degreeCelsius}，
即每十年 0.10 至 \qty{0.15}{\degreeCelsius}。
(b) 今天的全球速率，每十年 \qty{0.20}{\degreeCelsius}，为
格陵兰上空最快增温速度的 1.3 至 2 倍。(c) 单一地点的变化可能远大于
全球平均的变化，因为各地区的变化部分相互抵消：
应当把全球速率与全球速率相比。
\end{solution}

\begin{solution}{exo:g12:past-climates:15}
(a) 按每千年 \qty{12}{m} 计，\qty{14}{m} 需 $14/12 \approx
1.17$ 千年，约 \num{1170} 年：约 \num{9800} 年前。
(b) 否：海平面须至少比今天低 \qty{36}{m}，婆罗洲才能与亚洲相连；
一旦低于 \qty{-36}{m}，陆桥即被淹没，
婆罗洲重新成为岛屿。(c) $36/12 = 3$ 千年，约
\num{3000} 年；事实上，上升在约 \num{8200} 年前已放缓，
故实际所需时间更长。
\end{solution}

\begin{solution}{pb:g12:past-climates:1}
\textbf{1.} 冰本身借其中的重同位素记录
温度；冰中气泡内的空气记录二氧化碳。

\textbf{2.} $5.5 - (-10.6) = \qty{16.1}{\degreeCelsius}$。

\textbf{3.} $299 - 174 = \qty{125}{ppm}$。

\textbf{4.} 各点几乎落在一条直线上：两条记录同升同降，
但不能据此判断二者孰为因、孰为果。

\textbf{5.} 过去的记录以最古老的时间置于横轴的任一端绘制；
若读反方向，变暖便会被误读为变冷。

\textbf{6.} \omterm{def:g12:past-climates:glacial-cycle}{冰期}是冰河时代之内的寒冷时段，此时冰盖
远超其现今规模；\omterm{def:g12:past-climates:glacial-cycle}{间冰期}是两次\omterm{def:g12:past-climates:glacial-cycle}{冰期}之间的温暖时段，
此时冰盖接近今天的规模或
更小。

\textbf{7.} 处于\omterm{def:g12:past-climates:glacial-cycle}{间冰期}，始于约 \num{11000} 年前。

\textbf{8.} $\num{800000}/\num{98000} \approx 8$：约八个循环。

\textbf{9.} 轨道的形状。

\textbf{10.} 一个 \num{41000} 年的节律，即地轴倾角变化的
周期。

\textbf{11.} $-0.4 - (-9.4) = \qty{9.0}{\degreeCelsius}$。

\textbf{12.} $264 - 190 = \qty{74}{ppm}$。

\textbf{13.} \qty{6}{\degreeCelsius}；南极的变化 \qty{9.0}{\degreeCelsius}，
比全球变化大，约为全球变化的 1.5 倍。

\textbf{14.} $\num{18000} - \num{11000} = \num{7000}$ 年。

\textbf{15.} $6/7 \approx \qty{0.86}{\degreeCelsius}$ 每千年。

\textbf{16.} 一千年即一百个十年：$0.86/100 \approx
\qty{0.0086}{\degreeCelsius}$ 每十年。

\textbf{17.} $0.20 \times 100 = \qty{20}{\degreeCelsius}$ 每千年。

\textbf{18.} $20/0.86 \approx 23$ 倍。

\textbf{19.} $20/1.5 \approx 13$ 与 $20/1 = 20$：13 至 20 倍。

\textbf{20.} 今天的全球变暖速度\textbf{约快二十倍}，
比末次\omterm{def:g12:past-climates:glacial-cycle}{冰期}结束时的变暖快：是其平均速率的 23 倍，
为其最快速率的 13 至 20 倍。
\end{solution}
